% =====================================================================
%  ANEXO C — Casos de prueba funcionales
% =====================================================================
\chapter{CASOS DE PRUEBA FUNCIONALES}
\label{anx:casos}

Los casos de prueba se ejecutan sobre el \emph{development build} instalado
en un dispositivo Android, con una cuenta de prueba y el catálogo de máquinas
cargado en Firestore. En la columna \emph{Resultado} se registra si el caso
pasa o falla.

% TODO (autora): al ejecutar cada caso, reemplazar \cpres por "Pasa" o
% "Falla (ver observación)".
\newcommand{\cpres}{\textcolor{black!45}{\scriptsize\shortstack[l]{$\square$ Pasa\\$\square$ Falla}}}

{\footnotesize
\begin{longtable}{|L{1.15cm}|L{1.0cm}|L{5.3cm}|L{4.9cm}|C{1.35cm}|}
\caption{Casos de prueba funcionales}
\label{tab:casos-prueba}\\
\hline
\rowcolor{encabezado} \textbf{ID} & \textbf{RF} & \textbf{Pasos} & \textbf{Resultado esperado} &
\textbf{Resultado} \\
\hline
\endfirsthead
\hline
\rowcolor{encabezado} \textbf{ID} & \textbf{RF} & \textbf{Pasos} & \textbf{Resultado esperado} &
\textbf{Resultado} \\
\hline
\endhead
\endfoot
\multicolumn{5}{|l|}{\textbf{Autenticación, onboarding y perfil}}\\ \hline
CP-01 & 1.1 & Registrarse con nombre, apellido, correo válido y contraseña de
6 o más caracteres. & Se crea la cuenta y el perfil; se abre el
\emph{onboarding}. & \cpres \\ \hline
CP-02 & 1.1 & Registrarse con correo inválido, contraseñas distintas o de
menos de 6 caracteres. & Se muestra el mensaje de validación; no se crea la
cuenta. & \cpres \\ \hline
CP-03 & 1.2 & Iniciar sesión con correo y contraseña correctos, y luego con
una contraseña incorrecta. & El primero ingresa; el segundo muestra un mensaje
de error. & \cpres \\ \hline
CP-04 & 1.2 & Iniciar sesión con Google por primera vez; cerrar sesión y
volver a ingresar. & El perfil se crea solo la primera vez; no se
sobrescribe. & \cpres \\ \hline
CP-05 & 1.3 & Solicitar la recuperación de contraseña desde el inicio de
sesión. & Llega el correo de restablecimiento. & \cpres \\ \hline
CP-06 & 1.4 & Ingresar con un usuario sin perfil completo. & Se muestra el
\emph{onboarding}; las pestañas de la app no están disponibles. & \cpres \\
\hline
CP-07 & 1.4 & En datos físicos, ingresar edad 13 o estatura 20~cm. & Se
muestra el mensaje de rango; no se avanza. & \cpres \\ \hline
CP-08 & 1.4 & Completar objetivo, nivel y datos físicos y elegir ``Generar
con IA''. & Se guardan los datos; se abre la generación de rutina. & \cpres \\
\hline
CP-09 & 1.5 & Editar el objetivo y el peso en el perfil; cerrar sesión. & Los
cambios persisten; al cerrar sesión se vuelve a la bienvenida. & \cpres \\
\hline
\multicolumn{5}{|l|}{\textbf{Guía de máquinas y ejercicios}}\\ \hline
CP-10 & 2.1 & Abrir la pestaña Máquinas. & Se muestran las 35 máquinas en
cuadrícula. & \cpres \\ \hline
CP-11 & 2.1 & Buscar ``biceps'' y luego ``bíceps''. & Ambas búsquedas
devuelven los mismos resultados. & \cpres \\ \hline
CP-12 & 2.1 & Aplicar un filtro de grupo muscular y uno de tipo. & Solo se
muestran las máquinas que cumplen ambos filtros. & \cpres \\ \hline
CP-13 & 2.2 & Abrir la ficha de la prensa de piernas. & Se ven descripción,
músculos, guía de uso, errores, seguridad y variantes. & \cpres \\ \hline
CP-14 & 2.3 & En la ficha, revisar la sección de ejercicios. & Cada ejercicio
muestra su animación, músculo objetivo y dificultad. & \cpres \\ \hline
CP-15 & 2.2 & Revisar la galería de variantes de una ficha. & Se muestran
hasta ocho modelos comerciales. & \cpres \\ \hline
\multicolumn{5}{|l|}{\textbf{Escáner con realidad aumentada}}\\ \hline
CP-16 & 2.4 & Abrir el escáner y negar el permiso de la cámara. & Se muestra
el aviso del permiso; no se activa la cámara. & \cpres \\ \hline
CP-17 & 2.4 & Conceder el permiso y capturar una máquina. & Pasa por
``analizando''; si la confianza es de al menos 0,8 la máquina queda detectada, y
si no, se abre la lista con la sugerencia preseleccionada para confirmarla. &
\cpres \\
\hline
CP-18 & 2.4 & En la lista, buscar y elegir una máquina. & La ficha resumida
(nombre, tipo, músculos y ejercicios) se superpone sobre la cámara. & \cpres \\
\hline
CP-19 & 2.4 & Abrir el escáner desde una ficha; tocar ``escanear otra''. &
La máquina aparece directamente; luego se reinicia el escaneo. & \cpres \\
\hline
\multicolumn{5}{|l|}{\textbf{Rutinas y planificación}}\\ \hline
CP-20 & 3.1 & Con objetivo ``ganar masa'' y nivel principiante, generar una
rutina. & Rutina de 3 días (empuje, tirón, pierna), 4 series de 9
repeticiones, descanso 90~s. & \cpres \\ \hline
CP-21 & 3.1 & Con objetivo ``perder peso'' y nivel intermedio, generar una
rutina. & Rutina de 4 días con un día de cardio, 3 series de 15 repeticiones.
& \cpres \\ \hline
CP-22 & 3.2 & Crear una rutina personalizada de dos días con ejercicios. &
La rutina se guarda y queda activa. & \cpres \\ \hline
CP-23 & 3.2 & Intentar guardar sin nombre, o con un día sin ejercicios. & No
se guarda; se muestra el motivo. & \cpres \\ \hline
CP-24 & 3.3 & Desde una ficha, agregar un ejercicio a un día; repetir la
acción. & La primera vez se agrega; la segunda se informa el duplicado. &
\cpres \\ \hline
CP-25 & 3.3 & Activar otra rutina y eliminar una. & La activa se muestra en el
inicio; la eliminada desaparece. & \cpres \\ \hline
CP-26 & 3.3 & Tocar un ejercicio de la rutina asociado a una máquina. & Se abre
la ficha de la máquina. & \cpres \\ \hline
\multicolumn{5}{|l|}{\textbf{Nutrición: alimentos e IMC}}\\ \hline
CP-27 & 4.1 & Calcular el IMC con 70~kg y 175~cm. & IMC 22.9, categoría
normal. & \cpres \\ \hline
CP-28 & 4.1 & Calcular con peso 5~kg o estatura 20~cm. & Mensaje de
validación; no se calcula. & \cpres \\ \hline
CP-29 & 4.1 & Realizar un cálculo y revisar el historial. & El cálculo queda
registrado. & \cpres \\ \hline
CP-30 & 5.1 & Buscar ``pollo'', elegir un resultado e ingresar 150~g. & Se
listan alimentos; la vista previa y el registro muestran los mismos valores. &
\cpres \\ \hline
CP-31 & 5.2 & Tocar un alimento rápido (por ejemplo, ``Huevo''). & Se agrega
al día con sus calorías y macronutrientes. & \cpres \\ \hline
CP-32 & 5.2 & Ingresar un alimento manualmente con nombre y calorías. & Se
agrega al día. & \cpres \\ \hline
CP-33 & 5.3 & Revisar el resumen del día y eliminar un registro. & Las
calorías y macros se actualizan; el registro desaparece. & \cpres \\ \hline
\multicolumn{5}{|l|}{\textbf{Progreso, entrenamientos y reportes}}\\ \hline
CP-34 & 4.2 & En Progreso, registrar un peso (la estatura aparece
precargada). & Se crea el registro con su IMC. & \cpres \\ \hline
CP-35 & 4.2 & Registrar un peso de 0 o mayor de 500~kg. & Mensaje de
validación; no se guarda. & \cpres \\ \hline
CP-36 & 6.1 & Revisar el gráfico y alternar entre peso e IMC. & Muestra como
máximo los ocho últimos registros de la métrica elegida. & \cpres \\ \hline
CP-37 & 6.1 & Con registros de 80 y 77.5~kg, revisar la variación. & Se muestra
$-2.5$~kg; la categoría corresponde al último IMC. & \cpres \\ \hline
CP-38 & 6.2 & Marcar un día de la rutina como completado; marcarlo otra vez;
desmarcarlo. & Un solo registro por fecha; al desmarcar desaparece. & \cpres \\
\hline
CP-39 & 6.2 & Marcar entrenamientos en días consecutivos. & Racha, días de la
semana (lunes a domingo) y total se actualizan. & \cpres \\ \hline
CP-40 & 6.2 & Abrir el inicio con una rutina activa. & Muestra la rutina, sus
ejercicios, días por semana y registros de peso. & \cpres \\
\hline
\end{longtable}
}
